% Source: zcash/zips commit 2bdf39494320e5e04057ac93b1794031ebfb00b5, protocol/protocol.tex lines 10023-10076
% Spec version: v2026.7.0-344-g2bdf39 [NU6.2], section 5.4.1.11 "Equihash Generator"
% Pinned: 2026-10-24

\lsubsubsubsection{Equihash Generator}{equihashgen}

$\EquihashGen{n, k}$ is a specialized \hashFunction that maps an input
and an index to an output of length $n$ bits. It is used in \crossref{equihash}.

\newsavebox{\powtagbox}
\begin{lrbox}{\powtagbox}
\begin{bytefield}[bitwidth=0.16em]{128}
    \sbitbox{64}{64-bit $\ascii{ZcashPoW}$} &
    \sbitbox{32}{32-bit $n$} &
    \sbitbox{32}{32-bit $k$}
\end{bytefield}
\end{lrbox}

\newsavebox{\powcountbox}
\begin{lrbox}{\powcountbox}
\begin{bytefield}[bitwidth=0.16em]{32}
    \sbitbox{32}{32-bit $g$}
\end{bytefield}
\end{lrbox}

Let $\powtag := \Justthebox{\powtagbox}$.

Let $\powcount(g) := \Justthebox{\powcountbox}$.

\vspace{2ex}
\introlist
% Blech. Dijkstra was right \cite{EWD-831}.
Let $\EquihashGen{n, k}(S, i) := T_\barerange{h+1}{h+n}$, where
\begin{formulae}
  \item $m = \floor{\frac{512}{n}}$;
  \item $h = (i-1 \bmod m) \mult n$;
  \item $T = \BlakeTwobOf{(\mathnormal{n \mult m})}{\powtag,\, S \bconcat \powcount(\floor{\frac{i-1}{m}})}$.
\end{formulae}

Indices of bits in $T$ are 1-based.

$\BlakeTwobOf{\ell}{p, x}$ is defined in \crossref{concreteblake2}.

\securityrequirement{
$\BlakeTwobOf{\ell}{\powtag, x}$ must generate output that is sufficiently
unpredictable to avoid short-cuts to the \Equihash solution process.
It would suffice to model it as a \randomOracle.
}

\pnote{
When $\EquihashGen{}$ is evaluated for sequential indices, as
in the \Equihash solving process (\crossref{equihash}),
the number of calls to $\BlakeTwobGeneric$ can be reduced by a factor of
$\floor{\frac{512}{n}}$ in the best case (which is a factor of 2 for
$n = 200$).
}


